\textbf{1. Kiến thức cần nhớ}
\begin{itemize}
\item Cho hai đường tròn $(O; R)$ và $(O'; r)$ với $R \ge r$, gọi $d = OO'$ là đoạn nối tâm:
\begin{itemize}
\item Hai đường tròn \textbf{cắt nhau} $\iff \boxed{R - r < d < R + r}$ ($2$ điểm chung, đoạn nối tâm là đường trung trực của dây chung).
\item Hai đường tròn \textbf{tiếp xúc ngoài} $\iff \boxed{d = R + r}$ ($1$ tiếp điểm nằm giữa $O$ và $O'$).
\item Hai đường tròn \textbf{tiếp xúc trong} $\iff \boxed{d = R - r > 0}$ ($1$ tiếp điểm, $O, O'$ cùng phía đối với tiếp điểm).
\item Hai đường tròn \textbf{ở ngoài nhau} $\iff \boxed{d > R + r}$ (không giao nhau).
\item Đường tròn $(O)$ \textbf{đựng} $(O') \iff \boxed{d < R - r}$ (khi $d = 0$ gọi là hai đường tròn đồng tâm).
\end{itemize}
\end{itemize}

\textbf{2. Phương pháp giải các dạng}

%% DANG 9H5?17-1
\textbf{Dạng 1. Xác định vị trí tương đối của hai đường tròn}
\begin{enumerate}
\item Tính độ dài đoạn nối tâm $d = OO'$ và tính tổng $R + r$, hiệu $R - r$.
\item So sánh $d$ với $R + r$ và $R - r$ để kết luận vị trí tương đối.
\end{enumerate}

%% DANG 9H5?17-2
\textbf{Dạng 2. Tiếp tuyến chung, bài toán tổng hợp}
\begin{enumerate}
\item Vẽ tiếp tuyến chung: tiếp xúc ngoài có $3$ tiếp tuyến chung; ở ngoài nhau có $4$ tiếp tuyến chung; cắt nhau có $2$ tiếp tuyến chung ngoài; tiếp xúc trong có $1$ tiếp tuyến chung.
\item Đoạn nối tâm luôn là trục đối xứng của hình tạo bởi hai đường tròn.
\end{enumerate}
